\documentclass[11pt,a4paper]{article}

\usepackage[T1]{fontenc}
\usepackage[utf8]{inputenc}
\usepackage{lmodern}
\usepackage{amssymb,amsmath,amsthm,mathtools}
\usepackage[margin=2.5cm]{geometry}
\usepackage{parskip}
\usepackage[hidelinks]{hyperref}
\usepackage{microtype}

% Theorem environments (arXiv / journal style)
\theoremstyle{plain}
\newtheorem{theorem}{Theorem}[section]
\newtheorem{proposition}[theorem]{Proposition}
\newtheorem{lemma}[theorem]{Lemma}
\newtheorem{corollary}[theorem]{Corollary}

\theoremstyle{definition}
\newtheorem{definition}[theorem]{Definition}
\newtheorem{example}[theorem]{Example}
\newtheorem{remark}[theorem]{Remark}

% Notation
\newcommand{\Shape}{\mathrm{Shape}}
\newcommand{\Pos}{\mathrm{Pos}}
\newcommand{\Type}{\mathbf{Type}}
\newcommand{\id}{\mathrm{id}}
\newcommand{\src}{\mathrm{src}}
\newcommand{\tgt}{\mathrm{tgt}}
\newcommand{\rootpos}{\mathrm{root}}
\newcommand{\sub}{\mathrm{sub}}
\newcommand{\shift}{\mathrm{shift}}

\title{Compositional Tracking on Shipping Networks}
\author{Extracted from a machine-checked development}
\date{}

\begin{document}
\maketitle

\begin{abstract}
Given a directed graph~$G$ and a monoid~$M$, a \emph{compositional tracking}
of $M$-valued data on the shipping network of~$G$ assigns to every route an
element of~$M$, sending the empty route to the unit and sending concatenation
of routes to multiplication in~$M$. We prove that such trackings are in
bijection with arbitrary labellings of the edges of~$G$ by elements of~$M$:
restriction to single-edge routes and folding of edge labels along a path are
mutually inverse. Three monoid instances illustrate the construction---peak
temperature, cold-chain integrity, and licence accumulation---and a minimal
example shows that a single warm last-mile leg breaks an entire cold-chain
route. The proofs of all stated results have been machine-checked; we record
only the mathematical content.
\end{abstract}

\tableofcontents
\bigskip

%======================================================================
\section{Introduction}
%======================================================================

Shipping networks, presented as directed containers, organise nodes and
directed paths. Once that geometry is fixed, one often wishes to track
\emph{data along routes}: peak temperatures, cold-chain integrity flags,
accumulated licence requirements, and so on. Compositionality demands that the
value on a concatenated route be the monoid product of the values on the
factors, and that the empty route carry the monoid unit.

The present note isolates that structure. Working directly on the directed
container $\mathrm{network}(G)$ associated to a directed graph~$G$, a
compositional tracking with values in a monoid~$M$ is a family of maps from
positions to~$M$ satisfying a root law and a shift law. The main theorem is a
bijection: every tracking arises uniquely by folding an edge labelling along
paths, and every edge labelling extends uniquely to a tracking.

Categorically, a tracking is a functor from the free category on~$G$ to the
delooping of~$M$ (the one-object category with morphisms~$M$). The edge-folding
description is the elementary content of that functoriality.

The development works throughout in a single universe of types. Proofs are
indicated by the phrase ``by lean,'' meaning they have been checked by a proof
assistant rather than written out in prose.

\paragraph{What this note claims.}
In ordinary mathematical language: monoid-valued functors on the free category
of a directed graph are precisely edge labellings, extended by multiplication
along paths. There is no further coherence data beyond the monoid laws of~$M$
and the path algebra of~$G$.

%======================================================================
\section{Monoids}
%======================================================================

\begin{definition}[Monoid]\label{def:monoid}
A \textbf{monoid} $M$ consists of a type $\lvert M\rvert$, a binary operation
$\cdot : \lvert M\rvert \to \lvert M\rvert \to \lvert M\rvert$, and a unit
element $1 \in \lvert M\rvert$, such that for all $a,b,c \in \lvert M\rvert$,
\begin{equation}
  (a \cdot b) \cdot c \;=\; a \cdot (b \cdot c),
  \qquad
  1 \cdot a \;=\; a,
  \qquad
  a \cdot 1 \;=\; a.
\end{equation}
\end{definition}

We work with monoids as bundled structures rather than as typeclass instances.
When the monoid is clear we write $M$ both for the structure and for its
carrier.

%======================================================================
\section{Compositional trackings}
%======================================================================

Fix a directed graph~$G$ with node type~$V$ and edge type~$E$, and write
$\mathcal{N} := \mathrm{network}(G)$ for the associated directed container:
shapes are nodes; a position at~$v$ is a pair $\langle w,\rho \rangle$ of an
endpoint and a finite directed path $\rho$ from $v$ to~$w$; the root is the
empty path; sub returns the endpoint; and shift is path concatenation.

\begin{definition}[Compositional tracking]\label{def:tracking}
Let $M$ be a monoid. A \textbf{compositional tracking} of $M$-data on
$\mathcal{N}$ consists of a family
\begin{equation}
  d \;:\; \prod_{v \,:\, V}\;
    \Bigl(\Pos(\mathcal{N})(v) \;\longrightarrow\; M\Bigr)
\end{equation}
satisfying
\begin{enumerate}
  \item \textbf{(T-root)} for every node~$v$,
        \[
          d_v\bigl(\rootpos(v)\bigr) \;=\; 1;
        \]
  \item \textbf{(T-shift)} for every node~$v$, every position
        $\rho \in \Pos(\mathcal{N})(v)$, and every position
        $\sigma \in \Pos(\mathcal{N})\bigl(\sub(v,\rho)\bigr)$,
        \[
          d_v\bigl(\shift(v,\rho,\sigma)\bigr)
          \;=\;
          d_v(\rho) \cdot d_{\sub(v,\rho)}(\sigma).
        \]
\end{enumerate}
\end{definition}

Thus empty routes are labelled by the unit, and concatenated routes are labelled
by monoid products. Equivalently, $d$ is a functor $\mathcal{N}(G) \to BM$ from
the free category on~$G$ to the delooping of~$M$.

%======================================================================
\section{Folding along paths}
%======================================================================

\begin{definition}[Path fold]\label{def:fold}
Let $f : E \to M$ be a function labelling edges by monoid elements. The
\textbf{fold} of $f$ along a path is the map
\begin{equation}
  \mathrm{fold}_f \;:\;
  \prod_{v,w \,:\, V}\;
    \bigl(\mathrm{Path}_G(v,w) \;\longrightarrow\; M\bigr)
\end{equation}
defined by
\begin{align}
  \mathrm{fold}_f\bigl(\mathrm{nil}(v)\bigr) &\;=\; 1, \\
  \mathrm{fold}_f\bigl(\mathrm{cons}(e,p)\bigr)
  &\;=\;
  f(e) \cdot \mathrm{fold}_f(p).
\end{align}
\end{definition}

\begin{proposition}[Fold of the empty path]\label{prop:fold-nil}
For every edge labelling $f : E \to M$ and every node~$v$,
\[
  \mathrm{fold}_f\bigl(\mathrm{nil}(v)\bigr) \;=\; 1.
\]
\end{proposition}

\begin{proof}
By lean.
\end{proof}

\begin{proposition}[Fold preserves concatenation]\label{prop:fold-append}
For every edge labelling $f : E \to M$ and paths
$p : \mathrm{Path}_G(v,w)$, $q : \mathrm{Path}_G(w,u)$,
\[
  \mathrm{fold}_f\bigl(p \mathbin{{+}{+}} q\bigr)
  \;=\;
  \mathrm{fold}_f(p) \cdot \mathrm{fold}_f(q).
\]
\end{proposition}

\begin{proof}
By lean.
\end{proof}

The empty-path case is immediate from the left unit law of~$M$. The inductive
step for a cons-path uses associativity of multiplication.

%======================================================================
\section{From edge labellings to trackings and back}
%======================================================================

\begin{definition}[Tracking from an edge labelling]\label{def:of-edge}
Let $f : E \to M$. Define a tracking $\mathrm{ofEdge}(f)$ by
\begin{equation}
  d_v\bigl\langle w,\rho \bigr\rangle
  \;=\;
  \mathrm{fold}_f(\rho).
\end{equation}
The root law holds by Proposition~\ref{prop:fold-nil}; the shift law holds by
Proposition~\ref{prop:fold-append}.
\end{definition}

\begin{definition}[Single-edge position]\label{def:single-edge}
For an edge $e \in E$, the \textbf{single-edge position} at $\src(e)$ is
\begin{equation}
  \varepsilon_e
  \;=\;
  \bigl\langle
    \tgt(e),\;
    \mathrm{cons}\bigl(e,\, \mathrm{nil}(\tgt(e))\bigr)
  \bigr\rangle
  \;\in\;
  \Pos(\mathcal{N})\bigl(\src(e)\bigr).
\end{equation}
\end{definition}

\begin{definition}[Edge labelling from a tracking]\label{def:to-edge}
Let $t$ be a tracking with value family~$d$. Define
\begin{equation}
  \mathrm{toEdge}(t) \;:\; E \;\longrightarrow\; M,
  \qquad
  e \;\longmapsto\; d_{\src(e)}(\varepsilon_e).
\end{equation}
\end{definition}

\begin{lemma}[Fold recovers the tracking]\label{lem:fold-eq}
Let $t$ be a tracking with value family~$d$. For every path
$p : \mathrm{Path}_G(v,w)$,
\begin{equation}
  \mathrm{fold}_{\mathrm{toEdge}(t)}(p)
  \;=\;
  d_v\bigl\langle w,\, p \bigr\rangle.
\end{equation}
\end{lemma}

\begin{proof}
By lean.
\end{proof}

\begin{remark}[Idea of the argument]
The proof is by induction on~$p$. The nil case is the root law. For a cons-path
$\mathrm{cons}(e,p)$, the shift law applied to the single-edge position
$\varepsilon_e$ and the residual position $\langle w,p \rangle$ yields the
multiplicative step, after observing that shifting $\varepsilon_e$ by that
residual is definitionally $\langle w,\, \mathrm{cons}(e,p) \rangle$.
\end{remark}

\begin{theorem}[Round-trip: tracking to edges and back]\label{thm:of-to}
For every tracking~$t$,
\[
  \mathrm{ofEdge}\bigl(\mathrm{toEdge}(t)\bigr) \;=\; t.
\]
\end{theorem}

\begin{proof}
By lean.
\end{proof}

\begin{theorem}[Round-trip: edges to tracking and back]\label{thm:to-of}
For every edge labelling $f : E \to M$,
\[
  \mathrm{toEdge}\bigl(\mathrm{ofEdge}(f)\bigr) \;=\; f.
\]
\end{theorem}

\begin{proof}
By lean.
\end{proof}

\begin{theorem}[Tracking--edge bijection]\label{thm:tracking-bij}
For every directed graph~$G$ and monoid~$M$, there is a bijection
\begin{equation}
  \bigl\{\,\text{compositional trackings of $M$-data on
    $\mathrm{network}(G)$}\,\bigr\}
  \;\cong\;
  \bigl\{\,\text{functions } E \to M\,\bigr\},
\end{equation}
implemented by $\mathrm{toEdge}$ in one direction and $\mathrm{ofEdge}$ in the
other. The two constructions are mutually inverse by
Theorems~\ref{thm:of-to} and~\ref{thm:to-of}.
\end{theorem}

\begin{proof}
By lean.
\end{proof}

%======================================================================
\section{Monoid instances}
%======================================================================

\begin{example}[Peak temperature]\label{ex:maxtemp}
Let $\mathrm{MaxTemp} = \{\mathrm{Cold},\, \mathrm{Warm}\}$ with multiplication
\begin{equation}
  \mathrm{Cold} \cdot x \;=\; x,
  \qquad
  \mathrm{Warm} \cdot x \;=\; \mathrm{Warm},
\end{equation}
and unit $\mathrm{Cold}$. This is an idempotent monoid: a single warm leg along
a route dominates all cold legs, and subsequent refrigeration cannot undo a
prior breach of the temperature bound.
\end{example}

\begin{example}[Cold-chain integrity]\label{ex:coldchain}
Let $M = \{\mathsf{true},\, \mathsf{false}\}$ with multiplication given by
boolean conjunction and unit $\mathsf{true}$. Interpreting $\mathsf{true}$ as
``intact'' and $\mathsf{false}$ as ``broken,'' a route is intact if and only if
every leg is intact. Once a leg breaks the chain, no subsequent leg can repair
it.
\end{example}

\begin{example}[Licence accumulation]\label{ex:licences}
Let $L$ be a type of licence tokens, and take the monoid of finite lists of
elements of~$L$ under concatenation, with unit the empty list. Routes accumulate
the licences required by their individual legs. (Deduplication is not needed for
the algebraic structure.)
\end{example}

\begin{proposition}[Warm last-mile breaks the cold chain]\label{prop:warm-last-mile}
Consider the two-node digraph with nodes $\{\mathrm{Hub},\, \mathrm{Shop}\}$ and
a single edge $\mathrm{lastMile} : \mathrm{Hub} \to \mathrm{Shop}$, labelled
$\mathsf{false}$ in the cold-chain monoid of Example~\ref{ex:coldchain}. Let
$t$ be the tracking induced by this edge labelling. Then:
\begin{enumerate}
  \item the empty route at $\mathrm{Hub}$ is intact ($t$ returns
        $\mathsf{true}$);
  \item the single-edge route $\mathrm{Hub} \to \mathrm{Shop}$ is broken
        ($t$ returns $\mathsf{false}$).
\end{enumerate}
\end{proposition}

\begin{proof}
By lean.
\end{proof}

%======================================================================
\section{Discussion}
%======================================================================

\paragraph{What has been established.}
Compositional trackings on a shipping network are completely classified by
edge labellings with values in the target monoid. The reconstruction is
canonical: fold labels along paths in one direction, restrict to single-edge
routes in the other. The monoid laws of~$M$ and the path algebra of~$G$ are the
only ingredients.

\paragraph{Relation to functors into deloopings.}
The root and shift laws are exactly the requirements that $d$ preserve
identities and composition in the free category on~$G$. The edge-labelling
description is the usual fact that a functor out of a free category is
determined by its values on generating morphisms.

\paragraph{Scope.}
The development treats a single monoid of values and does not address
multi-parameter trackings, dependent value types, or interaction with
aggregation morphisms of networks. Those extensions sit on top of the same
bijection.

\paragraph{Machine checking.}
Every definition and theorem statement above corresponds to a declaration in a
machine-checked formalisation. The phrase ``by lean'' indicates that the proof
has been verified by a proof assistant.

%======================================================================
\section*{Acknowledgements}
%======================================================================

The mathematical content follows the construction of compositional trackings in
the supply-chain development that motivated this note. The presentation is a
prose rendering of a formal development intended for human review of the
theorem statements.

\begin{thebibliography}{9}

\bibitem{abbott2005}
M.~Abbott, T.~Altenkirch, and N.~Ghani.
\newblock Containers: constructing strictly positive types.
\newblock \emph{Theoretical Computer Science}, 342(1):3--27, 2005.

\bibitem{ahman2014}
D.~Ahman and T.~Uustalu.
\newblock Coalgebraic update lenses.
\newblock In \emph{MFPS}, volume~308 of \emph{ENTCS}, pages 25--48, 2014.

\end{thebibliography}

\end{document}
